\documentclass[10pt,twocolumn,a4paper]{article}
\usepackage[utf8]{inputenc}
\usepackage[T1]{fontenc}
\usepackage{lmodern}
\usepackage{microtype}
\usepackage[margin=2.2cm,columnsep=0.55cm]{geometry}
\usepackage{amsmath,amssymb,amsthm,mathtools}
\usepackage{bm}
\usepackage{physics}
\usepackage{graphicx}
\usepackage{booktabs}
\usepackage{array}
\usepackage{multirow}
\usepackage{tabularx}
\usepackage{enumitem}
\usepackage{float}
\usepackage{caption}
\usepackage{subcaption}
\usepackage{xcolor}
\usepackage{listings}
\usepackage{fancyvrb}
\usepackage{mdframed}
\usepackage{tikz}
\usetikzlibrary{arrows.meta,positioning,calc,shapes.geometric}
\usepackage{siunitx}
\usepackage[colorlinks=true,linkcolor=blue!60!black,
            citecolor=blue!60!black,urlcolor=blue!60!black]{hyperref}
\usepackage{cleveref}

\captionsetup{font=small,labelfont=bf,textfont=it}

% ─── Theorem environments ────────────────────────────────────────────
\newtheorem{theorem}{Theorem}[section]
\newtheorem{lemma}[theorem]{Lemma}
\newtheorem{proposition}[theorem]{Proposition}
\newtheorem{corollary}[theorem]{Corollary}
\theoremstyle{definition}
\newtheorem{definition}[theorem]{Definition}
\newtheorem{example}[theorem]{Example}
\theoremstyle{remark}
\newtheorem{remark}[theorem]{Remark}

% ─── BrutScript listing style ────────────────────────────────────────
\definecolor{kw}{RGB}{0,100,180}
\definecolor{cm}{RGB}{100,100,100}
\definecolor{str}{RGB}{160,32,32}
\definecolor{num}{RGB}{0,128,0}
\definecolor{bg}{RGB}{248,248,252}

\lstdefinelanguage{BrutScript}{
  morekeywords={source,signal,rate,layer,invert,derive,clamp,given,using,
                decompose,from,term,baseline,window,percentile,coord,
                regime,classify,when,failure,detect,otherwise,
                watch,and,or,not,emit,confidence,is,
                model,hub,task,input,output,
                explain,match,agree,diverge,trace,flag,annotate,with,
                first_match,aligns},
  sensitive=true,
  morecomment=[l]{--},
  morestring=[b]",
}

\lstset{
  language=BrutScript,
  basicstyle=\footnotesize\ttfamily,
  keywordstyle=\color{kw}\bfseries,
  commentstyle=\color{cm}\itshape,
  stringstyle=\color{str},
  numberstyle=\tiny\color{cm},
  backgroundcolor=\color{bg},
  frame=single,
  framerule=0.4pt,
  rulecolor=\color{cm!40},
  numbers=left,
  numbersep=5pt,
  xleftmargin=10pt,
  breaklines=true,
  showstringspaces=false,
  tabsize=2,
  captionpos=b,
}

% ─── Custom commands ─────────────────────────────────────────────────
\newcommand{\kB}{k_{\mathrm{B}}}
\newcommand{\BS}{\textsc{BrutScript}}
\newcommand{\Ees}{E_{\mathrm{es}}}
\newcommand{\Ea}{E_{\mathrm{a}}}
\newcommand{\EDV}{\mathrm{EDV}}
\newcommand{\ESV}{\mathrm{ESV}}
\newcommand{\EF}{\mathrm{EF}}
\newcommand{\CO}{\mathrm{CO}}
\newcommand{\SV}{\mathrm{SV}}
\newcommand{\HR}{\mathrm{HR}}
\newcommand{\SE}{S_{\mathrm{e}}}
\newcommand{\SK}{S_{\mathrm{k}}}
\newcommand{\ST}{S_{\mathrm{t}}}
\newcommand{\Rc}{R_c}
\newcommand{\PCHR}{\mathrm{PCHR}}
\newcommand{\CSCI}{\mathrm{CSCI}}
\newcommand{\TCC}{\mathrm{TCC}}
\newcommand{\Ra}{\mathrm{Ra}}
\newcommand{\Tr}{\mathrm{Tr}}
\newcommand{\vaso}{\eta}
\newcommand{\melan}{m}
\newcommand{\Hbconc}{[\mathrm{Hb}]_{\mathrm{tot}}}
\newcommand{\Soxy}{S_{\mathrm{O_2}}}

% ─────────────────────────────────────────────────────────────────────
\title{\textbf{BrutScript: A Partition-Theoretic Domain-Specific Language\\
for Auditable Physiological AI Pipelines}\\[0.3em]
\large Syntax, Semantics, and Formal Specification}

\author{
Kundai Farai Sachikonye\\[0.3em]
AIMe Registry for Artificial Intelligence\\
Experimental Bioinformatics, Maximus-von-Imhof-Forum 3\\
85354 Freising, Germany\\
\texttt{kundai.sachikonye@bitspark.com}
}

\date{\today}

\begin{document}

\twocolumn[
\begin{@twocolumnfalse}
\maketitle
\begin{abstract}
\noindent
We introduce \BS{}, a domain-specific language (DSL) for writing
\emph{auditable} physiological signal processing pipelines that can run
alongside black-box machine learning models and produce step-by-step
explanations of their outputs.  The problem \BS{} addresses is
fundamental: AI models trained on physiological sensor data are opaque
--- they map raw camera pixels or wrist-sensor streams to health labels
without exposing which physical quantities drove each prediction.  \BS{}
resolves this by providing a typed, traceable scripting vocabulary whose
primitives correspond one-to-one with the partition-theoretic quantities
derived in the Bounded Phase Space framework: Beer--Lambert optical
inversion layers, Partition-Coupled Heart Rate (PCHR) decomposition
terms, S-entropy health coordinates, Kuramoto regime classifiers,
Cross-Scale Coherence Index (CSCI) coupling pairs, and rambling--trembling
motor decomposition.  A \BS{} script declares signal sources, named
inversion layers, physiological decompositions, model bindings (including
HuggingFace Hub checkpoints), and alignment blocks that run both the
rule-based and neural pipelines simultaneously.  The compiler emits a
structured JSON execution trace in which every intermediate value is
named, time-stamped, and attributed to the source line that produced it.
Divergences between the rule-based trace and the neural model's
predictions are flagged with a pointer to the specific decomposition term
responsible.  We specify the complete grammar of \BS{}, define the
denotational semantics of each block type, and state the correctness
theorems that govern trace fidelity, type safety, and signal closure.
The language is implemented as a TypeScript compiler targeting the
BRUT Observatory web application, where it runs in real time against
live WebGPU rPPG data, face RGB inversions, and keyboard/mouse motor
sensors.

\vspace{4pt}
\noindent\textbf{Keywords:} domain-specific language, explainable AI,
physiological computing, remote photoplethysmography, S-entropy,
partition theory, cardiac equations of state, HuggingFace, audit trail,
wearable sensors.
\end{abstract}
\vspace{10pt}
\end{@twocolumnfalse}
]

% =====================================================================
\section{Introduction}
\label{sec:intro}
% =====================================================================

\subsection{The Black-Box Problem in Physiological AI}

Modern AI systems applied to physiological sensing are effective but
opaque.  A model trained on face video can output ``elevated stress'' at
14:32 without revealing which pixels drove the prediction, whether the
elevated heart rate was thermal, hypoxic, or autonomic in origin, whether
the cardiac oscillator was coherent or turbulent at the time, or whether
the model was operating in a data regime where its training generalises.
This opacity is not merely inconvenient --- it is a fundamental obstacle
to clinical adoption, regulatory approval, and developer debugging.

The standard response --- post-hoc explainability methods such as SHAP
or LIME --- applies input perturbations and observes output changes, but
this approach is agnostic to the biophysical meaning of the perturbed
features.  Knowing that the green channel had high SHAP weight says
nothing about whether the signal was carrying cardiac or vasodilatory
information, and says nothing about whether the layered optical model
would have agreed with the neural attribution.

\subsection{The Partition-Theoretic Alternative}

The Bounded Phase Space framework~\cite{sachikonye2024integration,
sachikonye2024fluid,sachikonye2024gas,sachikonye2024transport} provides
a first-principles mathematical description of every quantity that a
consumer physiological sensor touches.  Cardiac coherence $\Rc$ is a
Kuramoto order parameter~\cite{kuramoto1984}.  Heart rate decomposes into
intrinsic, metabolic, hypoxic, and autonomic components via the PCHR
equation~\cite{sachikonye2024disambiguation}.  RGB pixel values invert to
melanin index, haemoglobin concentration, oxygenation, and vasodilation
via layered Beer--Lambert optics~\cite{sachikonye2024optical}.  Keystroke
timing decomposes into rambling (supraspinal) and trembling (spinal)
motor components~\cite{sachikonye2024motor}.

Each of these quantities has a formal definition, a derivation from first
principles, and a computable form from sensor data.  \BS{} is a scripting
language whose vocabulary \emph{is} these quantities.  A \BS{} program
is a physics-grounded specification of what is happening in the pipeline,
and the compiler is the mechanism by which that specification becomes an
executable audit trail.

\subsection{Design Goals}

\BS{} is designed around four requirements:

\begin{enumerate}[leftmargin=*,nosep]
\item \textbf{Traceability.} Every intermediate value produced during
  execution must be named, time-stamped, and attributed to a source line.
\item \textbf{Physical fidelity.} The language primitives must correspond
  to physically meaningful quantities, not to ML feature indices or
  opaque tensors.
\item \textbf{Model alignment.} The language must be able to declare a
  neural model (including a HuggingFace Hub checkpoint) as a named
  participant in the pipeline and compare its output against the
  rule-based trace at each step.
\item \textbf{Real-time execution.} Scripts must compile to code that
  runs inside a WebGPU-accelerated browser at $\geq$1\,Hz against live
  sensor streams.
\end{enumerate}

\subsection{Contributions}

\begin{enumerate}[leftmargin=*,nosep]
\item A complete formal grammar for \BS{} (\cref{sec:grammar}).
\item Denotational semantics for each block type (\cref{sec:semantics}).
\item Correctness theorems governing type safety, trace fidelity, and
  signal closure (\cref{sec:theorems}).
\item A specification of the JSON trace format produced by the compiler
  (\cref{sec:trace}).
\item A reference vocabulary table mapping \BS{} identifiers to
  their partition-theoretic definitions (\cref{sec:vocabulary}).
\item An annotated implementation blueprint for the TypeScript
  compiler (\cref{sec:implementation}).
\end{enumerate}

% =====================================================================
\section{Background and Motivation}
\label{sec:background}
% =====================================================================

\subsection{Signal Provenance in rPPG Pipelines}

Remote photoplethysmography (rPPG) extracts cardiac signals from face
video~\cite{verkruysse2008,wang2017posrPPG,dehaan2013chrom}.  The BRUT
Observatory~\cite{sachikonye2024optical} extends this to a full
biophysical inversion: a WebGPU compute shader extracts per-patch cardiac
coherence $\Rc$ across a $64\times64$ face grid, the CPU layer extracts
a global BVP sample, and the algebraic inverter recovers the four-dimensional
optical state $(\melan, \Hbconc, \Soxy, \vaso)$ from mean face RGB at
1\,Hz.  From $\vaso$ the system derives skin temperature $T_{\mathrm{skin}}$
and the PCHR decomposition of heart rate into components $(\HR_{\mathrm{int}},
\Delta\HR_{\mathrm{met}}, \Delta\HR_{\mathrm{O}_2}, \Delta\HR_{\mathrm{auto}})$.

Every quantity in this chain has a mathematical definition
(\cref{tab:vocabulary}).  The problem is that when a neural model outputs
a label such as ``stress'' or ``cardiac risk'', its internal computation
does not respect these definitions --- it operates on raw feature vectors
and produces a scalar or categorical output without indicating which
physical pathway drove it.

\subsection{Sensor Disambiguation}

The sensor disambiguation framework~\cite{sachikonye2024disambiguation}
establishes formally (Theorem~1 therein) that a single sensor reading is
$k$-degenerate: $\HR = 90\,\mathrm{bpm}$ is consistent with exercise
recovery, fever, hypoxia, stress, heart failure, and altitude
acclimatisation simultaneously.  Disambiguation requires at least three
independent sensor projections plus two partition coupling equations.
The PCHR decomposition achieves this: given simultaneous $\HR$,
$T_{\mathrm{skin}}$, and $\Soxy$, the four components are uniquely
determined.  A neural model that does not perform this decomposition
internally cannot distinguish these physiological states.

\BS{} makes the PCHR decomposition a first-class language construct,
so a developer can write the decomposition explicitly alongside the
neural model's prediction and observe whether they agree.

\subsection{Motor Signals as Physiological Channels}

The neuromuscular derivation~\cite{sachikonye2024motor} establishes that
keystroke inter-keystroke intervals (IKI) are direct readouts of efferent
motor commands: the IKI sequence is a discretised version of the
continuous motor programme output.  The rambling--trembling decomposition
of cursor velocity~\cite{sachikonye2024posture} separates supraspinal
voluntary adjustments ($\Ra$, 0.05--0.5\,Hz) from spinal stabilisation
oscillations ($\Tr$, 0.5--3.0\,Hz).  The ratio $\Ra/(\Ra+\Tr)$
quantifies the balance between voluntary and reflex motor control.

\BS{} includes these motor channels as first-class source signals,
enabling scripts that combine cardiac, optical, and motor pathways
in a single audit trace.

\subsection{Cardiac Equations of State}

The cardiac equations of state (EOS)~\cite{sachikonye2024eos} provide a
closed-form algebraic inversion from BVP-derived features
$(\HR, A_{\mathrm{rel}}, T_{\mathrm{skin}}, \Soxy)$ to hemodynamic
state $(\Ees, \Ea, \EDV, \SV, \EF, \CO)$.  The end-systolic elastance
$\Ees$ quantifies left ventricular contractility; arterial elastance
$\Ea$ quantifies afterload; their ratio $\Ees/\Ea$ characterises
ventricular--arterial coupling efficiency.  These quantities are
derivable from consumer camera data without contact sensors.

\BS{} exposes the EOS inversion as a decompose block, so developers
can write scripts that reason about contractility and coupling efficiency
and compare these derived quantities against a neural model's output.

% =====================================================================
\section{Formal Grammar}
\label{sec:grammar}
% =====================================================================

\subsection{Lexical Structure}

\BS{} source files are UTF-8 text.  Comments begin with \texttt{--} and
extend to end of line.  Whitespace (space, tab, newline) is
non-significant except as a token separator.  String literals are
delimited by double quotes.  Numeric literals are IEEE-754 floating
point.  Identifiers match \texttt{[a-zA-Z\_][a-zA-Z0-9\_\.]*}.

\subsection{Context-Free Grammar}

The complete grammar is given in BNF notation below.  Non-terminals
are written in angle brackets; terminals are in \texttt{typewriter};
$\epsilon$ denotes the empty production; $\{X\}$ denotes zero or more
repetitions of $X$; $[X]$ denotes optional $X$.

\begin{mdframed}[backgroundcolor=bg,linecolor=cm!40,linewidth=0.4pt]
\scriptsize
\begin{verbatim}
<program>   ::= { <block> }

<block>     ::= <source-block>
              | <layer-block>
              | <decompose-block>
              | <watch-block>
              | <model-block>
              | <explain-block>

<source-block> ::=
  "source" IDENT "{" { <signal-decl> | <rate-decl> } "}"

<signal-decl> ::=
  "signal" IDENT { "," IDENT } [ "--" COMMENT ]

<rate-decl>   ::= "rate" NUM "hz"

<layer-block> ::=
  "layer" IDENT "from" <ident-list> "{"
    { <invert-stmt> | <derive-stmt> }
  "}"

<invert-stmt> ::=
  "invert" IDENT "from" IDENT
    [ "given" <ident-list> ]
    "using" IDENT
    [ "baseline" NUM ( "s" | "m" | "h" | "d" ) ]
    [ "sqrt_compress" ]
    [ "--" COMMENT ]

<derive-stmt> ::=
  "derive" IDENT "=" <expr>
    [ "clamp" "[" NUM "," NUM "]" ]
    [ "--" COMMENT ]

<decompose-block> ::=
  "decompose" IDENT "from" <ident-list> "{"
    { <term-stmt>
    | <coord-stmt>
    | <regime-stmt>
    | <failure-stmt>
    | <pair-stmt> }
  "}"

<term-stmt> ::=
  "term" IDENT "=" <expr> [ "--" COMMENT ]

<coord-stmt> ::=
  "coord" IDENT "=" <expr> [ "--" COMMENT ]

<regime-stmt> ::=
  "regime" "=" "classify" "(" IDENT ")" "{"
    { IDENT "when" IDENT <cmp-op> NUM }
    "otherwise" IDENT
  "}"

<failure-stmt> ::=
  "failure" "=" "detect" "{"
    { IDENT "when" <condition> }
    "none" "otherwise"
  "}"

<pair-stmt> ::=
  "pair" IDENT
    "predicted" "=" <expr>
    "observed"  "=" <expr>
    [ "--" COMMENT ]

<watch-block> ::=
  "watch" IDENT "{"
    "when" <condition>
    { ( "and" | "or" ) <condition> }
    "emit" STRING "confidence" NUM
    [ "--" COMMENT ]
  "}"

<condition> ::=
  IDENT <cmp-op> NUM
  | IDENT "is" IDENT
  | IDENT "aligns" IDENT
  | "not" <condition>

<cmp-op>    ::= ">" | "<" | ">=" | "<=" | "==" | "!="

<model-block> ::=
  "model" IDENT "{"
    "hub"    STRING
    "task"   IDENT
    "input"  "{" { <input-binding> } "}"
    "output" { IDENT }
  "}"

<input-binding> ::=
  IDENT { "," IDENT } "from" IDENT [ "." IDENT ]

<explain-block> ::=
  "explain" IDENT "{"
    "brut_prediction"  "=" <expr>
    "model_prediction" "=" IDENT "." IDENT
    "match" "{"
      "agree"   "when" <condition>
      "diverge" "when" <condition> "{"
        { <trace-stmt> | <flag-stmt> }
      "}"
    "}"
    [ "annotate" IDENT "with" "{" { <annot-item> } "}" ]
  "}"

<trace-stmt>  ::= "trace" IDENT [ "." IDENT ]
<flag-stmt>   ::= "flag" STRING
<annot-item>  ::= IDENT ":" IDENT

<ident-list>  ::= IDENT { "," IDENT }

<expr>        ::= <term> { ( "+" | "-" ) <term> }
<term>        ::= <factor> { ( "*" | "/" ) <factor> }
<factor>      ::= "(" <expr> ")"
               | IDENT [ "." IDENT ]
               | NUM
               | <builtin-call>
<builtin-call> ::= IDENT "(" [ <expr> { "," <expr> } ] ")"
\end{verbatim}
\end{mdframed}

\subsection{Reserved Identifiers}

The following identifiers are reserved as built-in functions
(\cref{tab:builtins}):

\begin{table}[H]
\centering
\caption{Built-in functions available in \BS{} expressions.}
\label{tab:builtins}
\footnotesize
\begin{tabularx}{\columnwidth}{lX}
\toprule
\textbf{Identifier} & \textbf{Semantics} \\
\midrule
\texttt{baseline(x,window,pct)} & Running $p$-th percentile of $x$ over rolling window \\
\texttt{shannon\_entropy(x)}    & $-\sum p_i \ln p_i$ over histogram of $x$ \\
\texttt{sqrt\_compress(x)}      & $\sqrt{\max(0.1,\, x)}$ (vasodilation stabiliser) \\
\texttt{mean(x)}                & Rolling mean \\
\texttt{std(x)}                 & Rolling standard deviation \\
\texttt{first\_match(...)}      & Returns first fired \texttt{watch} event \\
\texttt{classify(x)}            & Dispatch to enclosing regime block \\
\texttt{detect(...)}            & Dispatch to enclosing failure block \\
\texttt{log2(x)}                & $\log_2 x$ \\
\texttt{exp(x)}                 & $e^x$ \\
\texttt{cos(x)}                 & $\cos x$ \\
\texttt{abs(x)}                 & $|x|$ \\
\bottomrule
\end{tabularx}
\end{table}

% =====================================================================
\section{Denotational Semantics}
\label{sec:semantics}
% =====================================================================

We define the denotational semantics of each \BS{} block type by giving
its \emph{value domain}, \emph{evaluation rule}, and \emph{trace contribution}.

\subsection{Value Domains}

\begin{definition}[Value domains]
\label{def:domains}
Let $\mathbb{R}$ denote floating-point reals, $\mathbb{S}$ the set of
string literals, $\mathbb{B} = \{\top,\bot\}$ boolean truth values,
and $\mathbb{L}$ the set of regime labels.  A \emph{signal} is a
function $\sigma: \mathbb{T} \to \mathbb{R}$ where $\mathbb{T}$ is
discrete time indexed by the session clock in milliseconds.  The
\emph{environment} $\Gamma: \mathrm{Ident} \to \mathbb{R} \cup \mathbb{S}
\cup \mathbb{B} \cup \mathbb{L}$ maps identifiers to their current values.
\end{definition}

\subsection{Source Block}

\begin{definition}[Source semantics]
\label{def:source}
A \texttt{source} block with name $s$ and signal declarations
$\{n_1,\ldots,n_k\}$ at rate $f_s$ extends the environment by
\[
  \Gamma \leftarrow \Gamma \cup \{n_i \mapsto \sigma_{n_i}(t) \mid i=1,\ldots,k\}
\]
where $\sigma_{n_i}$ is the live sensor stream for channel $n_i$ sampled
at $f_s$ Hz.  The source block emits a trace entry
$\langle\texttt{source},\, s,\, t,\, \{n_i \mapsto v_i\}\rangle$
at each sample tick.
\end{definition}

\subsection{Layer Block}

\begin{definition}[Layer semantics]
\label{def:layer}
A \texttt{layer} block $\ell$ with invert statements
$\{\texttt{invert}\; y_j\; \texttt{from}\; x_j\; \texttt{given}\;
G_j\; \texttt{using}\; f_j\}_{j=1}^{m}$
evaluates each inversion in order:
\[
  y_j = f_j\!\left(x_j,\; \{g \mid g \in G_j\}\right)
\]
and extends the environment $\Gamma \leftarrow \Gamma \cup \{y_j \mapsto v_j\}$.
Each inversion emits a trace entry
$\langle\texttt{layer},\;\ell,\;\texttt{invert}\;y_j,\;t,\;\mathrm{inputs},\;v_j,\;\mathrm{line}\rangle$.
\end{definition}

The canonical inversion functions $f_j$ for optical inversion are given by:

\begin{align}
\melan &= -\frac{\ln \ell_B}{2\,\bar\mu_{\mathrm{mel},B} L_2}
\label{eq:melanin} \\
\Hbconc &= -\frac{\ln \tilde\ell_R}{2\,\bar\mu_{\Hb,R} L_{3,0}\,\vaso}
\label{eq:hbconc} \\
\Soxy &= \frac{\tilde\ell_G - t_{\mathrm{deox}}}{t_{\mathrm{oxy}} - t_{\mathrm{deox}}}
\label{eq:spo2} \\
\vaso &= \sqrt{\max\!\left(0.1,\; A^{\mathrm{ac}}_G / A^{\mathrm{base}}_G\right)}
\label{eq:vaso}
\end{align}
as derived in~\cite{sachikonye2024optical}
(equations~\ref{eq:melanin}--\ref{eq:vaso} here correspond to
equations~(5)--(8) therein).  The skin temperature follows from the
thermoregulatory linearisation
\begin{equation}
T_{\mathrm{skin}} = 33 + 4(\vaso - 1)\ {}^{\circ}\mathrm{C},
\quad T_{\mathrm{skin}} \in [27,37]
\label{eq:tskin}
\end{equation}
as established in~\cite{sachikonye2024optical} and
\cite{sachikonye2024disambiguation}.

\subsection{Decompose Block}

\begin{definition}[Decompose semantics]
\label{def:decompose}
A \texttt{decompose} block $d$ with term statements
$\{\texttt{term}\; z_i = e_i\}_{i=1}^{n}$ evaluates each expression
$e_i$ in the current environment and binds $z_i \leftarrow \llbracket e_i
\rrbracket_\Gamma$, extending $\Gamma$ after each binding so later terms
may reference earlier ones.  Each term emits a trace entry
$\langle\texttt{decompose},\;d,\;\texttt{term}\;z_i,\;t,\;v_i,\;\mathrm{line}\rangle$.
\end{definition}

The PCHR decomposition~\cite{sachikonye2024disambiguation} is the
canonical instance, with terms:
\begin{align}
\HR_{\mathrm{int}} &= \mathrm{baseline}(\HR, 7\,\mathrm{d}, p_5)
\label{eq:hr_int} \\
\Delta\HR_{\mathrm{met}} &= \alpha_T \cdot (T_{\mathrm{skin}} - T_{\mathrm{rest}}) \cdot \HR_{\mathrm{int}}
\label{eq:dhr_met} \\
\Delta\HR_{\mathrm{O}_2} &= \beta_{O_2} \cdot (1 - \Soxy) \cdot \HR_{\mathrm{int}}
\label{eq:dhr_o2} \\
\Delta\HR_{\mathrm{auto}} &= \HR_{\mathrm{obs}} - \HR_{\mathrm{int}} - \Delta\HR_{\mathrm{met}} - \Delta\HR_{\mathrm{O}_2}
\label{eq:dhr_auto}
\end{align}
with coupling constants $\alpha_T = 0.08\,{}^{\circ}\mathrm{C}^{-1}$ (from
$Q_{10}=2.3$)~\cite{sachikonye2024disambiguation,sachikonye2024optical}
and $\beta_{O_2} = 0.15$ per 10\% SpO$_2$ drop.

The S-entropy health coordinates~\cite{sachikonye2024disambiguation} are:
\begin{align}
\SK &= \frac{\mathrm{RMSSD} \cdot (60000/\HR)}{\SK_{\max}}
\label{eq:sk} \\
\ST &= 1 - e^{-t_{\mathrm{awake}}/24\mathrm{h}} \cos(2\pi t_{\mathrm{awake}}/24\mathrm{h})
\label{eq:st} \\
\SE &= -\sum_i p_i \log_2 p_i \;/\; \log_2 N_{\mathrm{bins}}
\label{eq:se}
\end{align}
The healthy region is $\mathcal{H} = [0.4,0.9]\times[0.2,0.8]\times[0.6,0.95]
\subset [0,1]^3$~\cite{sachikonye2024disambiguation}.

The CSCI~\cite{sachikonye2024disambiguation} is:
\begin{equation}
\CSCI = 1 - \frac{1}{N_p}\sum_{(i,j)\in\mathcal{P}}
\left|\frac{r_{ij}^{\mathrm{obs}} - r_{ij}^{\mathrm{pred}}}{r_{ij}^{\mathrm{pred}}}\right|
\label{eq:csci}
\end{equation}
with the partition-predicted coupling pairs listed in~\cref{tab:coupling}.

\begin{table}[H]
\centering
\caption{Partition-predicted coupling pairs used in CSCI computation.}
\label{tab:coupling}
\footnotesize
\begin{tabularx}{\columnwidth}{lX}
\toprule
\textbf{Pair} & \textbf{Predicted ratio} \\
\midrule
$\Rc$ vs $\HR$ & $\partial\Rc/\partial\HR = -4\pi^2\mathrm{CV}^2\Rc/\HR$ \\
$\HR$ vs $\Delta T$ & $\Delta\HR/\Delta T = \alpha_T \HR_{\mathrm{int}}$ \\
$\HR$ vs $\Soxy$ & $\Delta\HR/\Delta\Soxy = -\beta_{O_2}\HR_{\mathrm{int}}/10$ \\
$\mathrm{RMSSD}$ vs $\HR$ & $\mathrm{RMSSD} \propto \HR^{-1.5}$ \\
\bottomrule
\end{tabularx}
\end{table}

\subsection{Watch Block}

\begin{definition}[Watch semantics]
\label{def:watch}
A \texttt{watch} block $w$ with compound condition
$C = c_1 \wedge c_2 \wedge \cdots \wedge c_n$ evaluates each
sub-condition $c_i \in \mathbb{B}$ in the current environment.
If $\bigwedge_i c_i = \top$, the watch block emits
$\langle\texttt{fired},\;w,\;t,\;\mathrm{label},\;\mathrm{confidence}\rangle$
together with a per-condition trace:
$\langle c_i,\;\llbracket c_i \rrbracket_\Gamma,\;\mathrm{value},\;\mathrm{line}\rangle$
for each $i$.  If any $c_i = \bot$, the block emits
$\langle\texttt{not\_fired},\;w,\;t,\;\{c_i \mapsto \bot : c_i = \bot\}\rangle$.
\end{definition}

This per-condition tracing is the primary mechanism by which \BS{} makes
the reasoning of the rule-based pipeline auditable: the developer can
read exactly which sub-condition was true, what its measured value was,
and which source line declared it.

\subsection{Model Block}

\begin{definition}[Model semantics]
\label{def:model}
A \texttt{model} block $m$ with hub identifier $h$, task $\tau$, and
input bindings $\{(q_j, s_j)\}$ declares a deferred inference call.
At evaluation time, the block assembles an input tensor
$\mathbf{x} = \{q_j \mapsto \Gamma(s_j)\}$, submits it to the
HuggingFace Inference API endpoint resolved from $h$, and binds
each declared output identifier $o_k$ to the corresponding tensor
slice of the model response.  Each model block emits a trace entry
$\langle\texttt{model},\;m,\;t,\;\mathbf{x},\;\{o_k \mapsto v_k\}\rangle$.
\end{definition}

The model block is not a training construct --- it invokes a
\emph{pre-trained} checkpoint at inference time.  The input bindings
use identifiers that must exist in $\Gamma$, enforcing that the model
receives values produced by prior \texttt{source}, \texttt{layer}, and
\texttt{decompose} blocks.  This makes the model's inputs auditable:
the trace shows exactly which decomposition terms were passed as features.

\subsection{Explain Block}

\begin{definition}[Explain semantics]
\label{def:explain}
An \texttt{explain} block $e$ evaluates:
\begin{enumerate}[nosep]
\item $b \leftarrow \llbracket \texttt{brut\_prediction} \rrbracket_\Gamma$
  (the rule-based event or label).
\item $m \leftarrow \llbracket \texttt{model\_prediction} \rrbracket_\Gamma$
  (the neural model's output).
\item The \texttt{agree} condition: if $\llbracket \texttt{agree.when}
  \rrbracket_\Gamma = \top$, emit
  $\langle\texttt{agree},\;e,\;t,\;b,\;m\rangle$.
\item The \texttt{diverge} condition: if $\llbracket \texttt{diverge.when}
  \rrbracket_\Gamma = \top$, evaluate all \texttt{trace} and \texttt{flag}
  statements and emit $\langle\texttt{diverge},\;e,\;t,\;b,\;m,\;
  \mathrm{traces},\;\mathrm{flags}\rangle$.
\item If \texttt{annotate} is present, attach the named decomposition
  blocks to the model prediction entry in the trace.
\end{enumerate}
\end{definition}

% =====================================================================
\section{Correctness Theorems}
\label{sec:theorems}
% =====================================================================

\begin{theorem}[Type safety]
\label{thm:type}
Let $\mathcal{P}$ be a well-formed \BS{} program (i.e., one that parses
without error under the grammar of \cref{sec:grammar}).  If every
identifier referenced in an expression $e$ is bound in the environment
$\Gamma$ at the point of evaluation, then $\llbracket e \rrbracket_\Gamma
\in \mathbb{R}$ (or the appropriate declared domain).
\end{theorem}

\begin{proof}[Proof sketch]
By structural induction on the grammar.  Base case: numeric literals are
in $\mathbb{R}$; identifiers in $\Gamma$ are bound by a prior \texttt{source},
\texttt{layer}, or \texttt{decompose} block.  Inductive step: each
arithmetic operator and built-in function maps $\mathbb{R}^k \to \mathbb{R}$.
Regime and failure classifiers map $\mathbb{R} \to \mathbb{L}$.
Watch conditions map $\mathbb{R} \times \mathbb{R} \to \mathbb{B}$.
No block can reference an identifier bound in a later block (programs are
evaluated sequentially), so the environment is always well-defined at the
point of use. $\square$
\end{proof}

\begin{theorem}[Signal closure]
\label{thm:closure}
A well-formed \BS{} program is \emph{signal-closed} if every identifier
referenced in a \texttt{layer}, \texttt{decompose}, \texttt{watch},
\texttt{model}, or \texttt{explain} block is either (a) declared in a
\texttt{source} block that appears before it in program order, (b) bound
by a prior \texttt{layer} or \texttt{decompose} statement, or (c) a
built-in function.  Signal closure is decidable by a single forward pass
over the program.
\end{theorem}

\begin{proof}
The compiler maintains a scope set $\mathcal{S}$ initialised to the
built-in function names.  Processing each block in order: a
\texttt{source} block adds its declared signal names to $\mathcal{S}$;
a \texttt{layer} or \texttt{decompose} statement adds its output name
after verifying all input names are in $\mathcal{S}$; all other
references are checked against $\mathcal{S}$ and rejected if absent.
The single-pass algorithm runs in $O(|\mathcal{P}|)$ time.  $\square$
\end{proof}

\begin{theorem}[Trace fidelity]
\label{thm:fidelity}
Let $\mathcal{P}$ be a signal-closed \BS{} program and let
$\mathcal{T}_\mathcal{P}(t)$ be the trace produced at time $t$.  Then
for every intermediate value $v$ produced during evaluation,
$\mathcal{T}_\mathcal{P}(t)$ contains a unique entry
$\langle\mathrm{block},\;\mathrm{step},\;t,\;\mathrm{inputs},\;v,\;\mathrm{line}\rangle$
where $\mathrm{line}$ identifies the source statement that produced $v$.
\end{theorem}

\begin{proof}
By inspection of the evaluation rules in \cref{sec:semantics}: each
evaluation rule emits exactly one trace entry per bound identifier per
time step.  The \texttt{line} field is populated by the parser, which
attaches source position to each AST node.  Uniqueness follows from
sequential evaluation within a block.  $\square$
\end{proof}

\begin{theorem}[Divergence completeness]
\label{thm:divergence}
If a \BS{} \texttt{explain} block fires its \texttt{diverge} branch,
then for every \texttt{trace} statement listed in the diverge body,
the trace $\mathcal{T}_\mathcal{P}(t)$ contains the complete evaluation
history of the referenced decomposition term, including all its inputs
and the source lines that produced them.
\end{theorem}

\begin{proof}
Follows from \cref{thm:fidelity}: every term in a \texttt{decompose}
block has a trace entry by the time the \texttt{explain} block evaluates.
The \texttt{trace} statement retrieves this entry by identifier lookup
in $\mathcal{T}_\mathcal{P}(t)$ and includes it in the diverge output.
Since the explain block is evaluated after all referenced blocks (signal
closure), no entry is missing.  $\square$
\end{proof}

% =====================================================================
\section{Execution Trace Format}
\label{sec:trace}
% =====================================================================

The compiler emits traces as newline-delimited JSON (NDJSON), one
object per event.  \cref{lst:trace} shows a representative trace
fragment for a stress-onset audit at time $t = 14{:}32{:}07$.

\begin{lstlisting}[
  language={},
  caption={Representative BrutScript execution trace (NDJSON).},
  label={lst:trace},
  basicstyle=\scriptsize\ttfamily,
  numbers=none
]
{"t":87127000,"block":"skin_optics",
 "step":"invert melanin","line":7,
 "in":{"b":0.41},"out":{"melanin":0.18}}

{"t":87127000,"block":"skin_optics",
 "step":"invert hb_conc","line":8,
 "in":{"r":0.38,"melanin":0.18},
 "out":{"hb_conc":13.2}}

{"t":87127000,"block":"skin_optics",
 "step":"derive t_skin","line":11,
 "in":{"vasodilation":1.02},
 "out":{"t_skin":33.08}}

{"t":87127000,"block":"pchr",
 "step":"term dhr_metabolic","line":16,
 "in":{"t_skin":33.08,"hr_int":58.0},
 "out":{"dhr_metabolic":0.37}}

{"t":87127000,"block":"pchr",
 "step":"term dhr_autonomic","line":18,
 "in":{"hr":97,"dhr_met":0.37,"dhr_o2":0.0},
 "out":{"dhr_autonomic":18.6}}

{"t":87127000,"block":"s_entropy",
 "step":"regime classify","line":23,
 "in":{"rc_mean":0.71},"out":{"regime":"turbulent"}}

{"t":87127000,"block":"watch stress_onset",
 "step":"cond dhr_autonomic > 15","line":38,
 "eval":true,"value":18.6}

{"t":87127000,"block":"watch stress_onset",
 "step":"cond regime is turbulent","line":39,
 "eval":true,"value":"turbulent"}

{"t":87127000,"block":"watch stress_onset",
 "step":"cond rt_ratio < 0.3","line":40,
 "eval":true,"value":0.21}

{"t":87127000,"block":"watch stress_onset",
 "step":"FIRED","emit":"stress_onset",
 "confidence":0.8,"line":41}

{"t":87127000,"block":"model stress_detector",
 "step":"inference","line":45,
 "in":{"dhr_auto":18.6,"rt_ratio":0.21,
       "mean_iki":142.0},
 "out":{"stress_score":0.34}}

{"t":87127000,"block":"explain stress_audit",
 "step":"diverge","brut":"stress_onset",
 "model_score":0.34,"line":50,
 "flag":"model may be confounding thermal
         with autonomic drive",
 "traces":{"pchr.dhr_metabolic":0.37,
           "skin_optics.t_skin":33.08}}
\end{lstlisting}

Every trace entry carries: a session-relative timestamp \texttt{t} in
milliseconds, the block name and step name, the source line number,
the input values, and the output value(s).  Watch entries additionally
carry the boolean evaluation result and the measured value.  Explain
entries carry the BRUT prediction, the model output, and the full text
of any flags.

% =====================================================================
\section{Reference Vocabulary}
\label{sec:vocabulary}
% =====================================================================

\cref{tab:vocabulary} provides the complete mapping from \BS{}
identifiers to their partition-theoretic definitions.  Every identifier
in the table is a legal \BS{} signal name, term name, or output name
in the corresponding block type.

\begin{table*}[ht]
\centering
\caption{BrutScript reference vocabulary. Column ``Block'' indicates
the block type where the identifier is produced or consumed.
Column ``Definition'' cites the equation or paper where the quantity
is formally defined.}
\label{tab:vocabulary}
\footnotesize
\begin{tabularx}{\textwidth}{llllX}
\toprule
\textbf{Identifier} & \textbf{Type} & \textbf{Units} & \textbf{Block} & \textbf{Definition / provenance} \\
\midrule
\multicolumn{5}{l}{\textit{rPPG source signals}} \\
\texttt{bvp}            & $\mathbb{R}$ & a.u.      & source & Global BVP sample from WebGPU rPPG shader~\cite{sachikonye2024optical} \\
\texttt{rc\_mean}       & $[0,1]$      & —         & source & Spatial mean Kuramoto coherence $\Rc$~\cite{kuramoto1984,sachikonye2024disambiguation} \\
\texttt{rc\_std}        & $[0,1]$      & —         & source & Spatial std of $\Rc$ (heterogeneity)~\cite{sachikonye2024optical} \\
\texttt{snr}            & $\mathbb{R}$ & dB        & source & Mean SNR over face patches~\cite{sachikonye2024optical} \\
\texttt{hr}             & $\mathbb{R}$ & bpm       & source & Heart rate from BVP analyser~\cite{sachikonye2024optical} \\
\texttt{rmssd}          & $\mathbb{R}$ & ms        & source & RMSSD of RR intervals~\cite{malik1996,sachikonye2024disambiguation} \\
\texttt{sk}             & $[0,1]$      & —         & source & S-entropy knowledge coordinate $\SK$~\cite{sachikonye2024disambiguation} \\
\texttt{st}             & $[0,1]$      & —         & source & S-entropy temporal coordinate $\ST$~\cite{sachikonye2024disambiguation} \\
\texttt{se}             & $[0,1]$      & —         & source & S-entropy entropy utilisation $\SE$~\cite{sachikonye2024disambiguation} \\
\midrule
\multicolumn{5}{l}{\textit{Face RGB source signals}} \\
\texttt{r, g, b}        & $[0,1]$      & —         & source & Mean normalised RGB from face ROI~\cite{sachikonye2024optical} \\
\texttt{bvp.ac\_amplitude} & $\mathbb{R}$ & a.u.   & source & AC pulsatile amplitude of BVP~\cite{sachikonye2024optical} \\
\midrule
\multicolumn{5}{l}{\textit{Motor source signals}} \\
\texttt{mean\_iki}      & $\mathbb{R}$ & ms        & source & Mean inter-keystroke interval~\cite{sachikonye2024motor} \\
\texttt{rt\_ratio}      & $[0,1]$      & —         & source & Rambling power / (Rambling + Trembling)~\cite{sachikonye2024posture} \\
\texttt{blinks\_min}    & $\mathbb{R}$ & min$^{-1}$& source & Blinks per minute from MediaPipe~\cite{sachikonye2024optical} \\
\midrule
\multicolumn{5}{l}{\textit{Optical inversion layer outputs}} \\
\texttt{melanin}        & $[0,1]$      & —         & layer  & Melanin index $\melan$ via blue-channel inversion~\cite{sachikonye2024optical}, \cref{eq:melanin} \\
\texttt{hb\_conc}       & $\mathbb{R}$ & g/dL      & layer  & Total haemoglobin $\Hbconc$ via red-channel inversion~\cite{sachikonye2024optical}, \cref{eq:hbconc} \\
\texttt{spo2}           & $[0,1]$      & —         & layer  & Oxygenation $\Soxy$ via green-channel inversion~\cite{sachikonye2024optical}, \cref{eq:spo2} \\
\texttt{vasodilation}   & $[0.6,1.6]$  & —         & layer  & Vasodilation factor $\vaso$ from AC amplitude~\cite{sachikonye2024optical}, \cref{eq:vaso} \\
\texttt{t\_skin}        & $[27,37]$    & ${}^\circ$C & layer & Skin temperature via linearisation~\cite{sachikonye2024optical,sachikonye2024disambiguation}, \cref{eq:tskin} \\
\midrule
\multicolumn{5}{l}{\textit{PCHR decomposition terms}} \\
\texttt{hr\_intrinsic}  & $\mathbb{R}$ & bpm       & decompose & Sino-atrial intrinsic rate (7-day $p_5$ baseline)~\cite{sachikonye2024disambiguation}, \cref{eq:hr_int} \\
\texttt{dhr\_metabolic} & $\mathbb{R}$ & bpm       & decompose & Thermal drive $\Delta\HR_{\mathrm{met}}$~\cite{sachikonye2024disambiguation,sachikonye2024optical}, \cref{eq:dhr_met} \\
\texttt{dhr\_hypoxic}   & $\mathbb{R}$ & bpm       & decompose & Hypoxic compensation $\Delta\HR_{O_2}$~\cite{sachikonye2024disambiguation}, \cref{eq:dhr_o2} \\
\texttt{dhr\_autonomic} & $\mathbb{R}$ & bpm       & decompose & Genuine neural drive $\Delta\HR_{\mathrm{auto}}$~\cite{sachikonye2024disambiguation}, \cref{eq:dhr_auto} \\
\midrule
\multicolumn{5}{l}{\textit{S-entropy and regime}} \\
\texttt{sk, st, se}     & $[0,1]^3$    & —         & decompose & S-entropy health coordinates~\cite{sachikonye2024disambiguation}, \cref{eq:sk,eq:st,eq:se} \\
\texttt{regime}         & $\mathbb{L}$ & —         & decompose & Kuramoto regime: phase\_locked $|$ coherent $|$ cascade $|$ aperture $|$ turbulent~\cite{sachikonye2024disambiguation} \\
\texttt{failure}        & $\mathbb{L}$ & —         & decompose & Failure mode: rigidity $|$ decoherence $|$ none~\cite{sachikonye2024disambiguation} \\
\texttt{csci}           & $[0,1]$      & —         & decompose & Cross-Scale Coherence Index~\cite{sachikonye2024disambiguation}, \cref{eq:csci} \\
\midrule
\multicolumn{5}{l}{\textit{Cardiac EOS outputs}} \\
\texttt{ees}            & $\mathbb{R}$ & mmHg/mL   & decompose & End-systolic elastance (contractility)~\cite{sachikonye2024eos} \\
\texttt{ea}             & $\mathbb{R}$ & mmHg/mL   & decompose & Arterial elastance (afterload)~\cite{sachikonye2024eos} \\
\texttt{ef}             & $[0,1]$      & —         & decompose & Ejection fraction $\EF = \SV/\EDV$~\cite{sachikonye2024eos} \\
\texttt{co}             & $\mathbb{R}$ & L/min     & decompose & Cardiac output~\cite{sachikonye2024eos,sachikonye2024cardiovascular} \\
\texttt{ees\_ea\_ratio} & $\mathbb{R}$ & —         & decompose & V-A coupling ratio $\Ees/\Ea$~\cite{sachikonye2024eos} \\
\bottomrule
\end{tabularx}
\end{table*}

% =====================================================================
\section{Complete Example Script}
\label{sec:example}
% =====================================================================

\cref{lst:full} presents a complete \BS{} script auditing a
cardiac stress episode.  The script is the formal basis for the
implementation blueprint in \cref{sec:implementation}.

\begin{lstlisting}[
  language=BrutScript,
  caption={Complete BrutScript stress audit script.
    All identifiers map to entries in \cref{tab:vocabulary}.},
  label={lst:full},
  float=*
]
-- BrutScript: stress episode audit
-- All primitives ground to partition theory (see vocabulary table)

source rppg {
  signal bvp, rc_mean, rc_std, snr
  signal hr, rmssd, sk, st, se
  rate   30hz
}

source face_rgb {
  signal r, g, b        -- mean normalised RGB from face ROI
  rate   1hz
}

source motor {
  signal mean_iki       -- inter-keystroke interval (ms)
  signal rt_ratio       -- rambling / (rambling + trembling)
  signal blinks_min
  rate   1hz
}

-- Stage 1: Beer-Lambert optical inversion (layered-optical-ppg)
layer skin_optics from face_rgb {
  invert melanin      from b  using beer_lambert.blue
  -- eq. (5): m = -ln(l_B) / (2 * mu_mel_B * L2)

  invert hb_conc      from r  given melanin
                              using beer_lambert.red
  -- eq. (6): [Hb] = -ln(l_R_tilde) / (2 * mu_Hb_R * L3 * eta)

  invert spo2         from g  given melanin, hb_conc
                              using beer_lambert.green
  -- eq. (7): S_O2 = (l_G_tilde - t_deox) / (t_oxy - t_deox)

  invert vasodilation from bvp.ac_amplitude
                              baseline 30s  sqrt_compress
  -- eq. (8): eta = sqrt(max(0.1, A_ac / A_base))

  derive t_skin = 33.0 + 4.0 * (vasodilation - 1.0)  clamp [27, 37]
  -- eq. (9): T_skin = 33 + 4*(eta-1), thermoregulatory linearisation
}

-- Stage 2: PCHR decomposition (sensor-disambiguation)
decompose pchr from rppg, skin_optics {
  term hr_intrinsic  = baseline(hr, window 7d, percentile p5)
  -- 7-day 5th-percentile deep-sleep HR = sino-atrial intrinsic rate

  term dhr_metabolic = 0.08 * (t_skin - 33.0) * hr_intrinsic
  -- alpha_T = 0.08 /degC from Q10 = 2.3 (Theorem alpha_T)

  term dhr_hypoxic   = 0.15 * (1.0 - spo2) * hr_intrinsic
  -- beta_O2 = 0.15 per 10% SpO2 drop (Theorem beta_O2)

  term dhr_autonomic = hr - hr_intrinsic - dhr_metabolic - dhr_hypoxic
  -- residual = genuine sympathovagal drive (eq. dHR_auto)
}

-- Stage 3: S-entropy and regime (sensor-disambiguation)
decompose s_entropy from rppg {
  coord sk = rmssd * (60000.0 / hr) / 1000.0
  coord se = shannon_entropy(rr_intervals) / log2(32.0)

  regime = classify(rc_mean) {
    phase_locked  when rc_mean >= 0.947
    coherent      when rc_mean >= 0.930
    cascade       when rc_mean >= 0.900
    aperture      when rc_mean >= 0.850
    turbulent     otherwise
  }

  failure = detect {
    rigidity    when rc_mean > 0.95 and se < 0.5
    decoherence when rc_mean < 0.30
    none        otherwise
  }
}

-- Stage 4: Cross-scale coupling index (sensor-disambiguation, eq. CSCI)
decompose coupling from rppg, skin_optics {
  pair hr_vs_temp
    predicted = 0.08 * hr_intrinsic
    observed  = dhr_metabolic / (t_skin - 33.0 + 0.001)

  pair rmssd_vs_hr
    predicted = hr ^ -1.5
    observed  = rmssd / 1000.0

  -- csci = 1 - mean(|obs - pred| / |pred|)
  term csci = 1.0 - mean(abs(hr_vs_temp.error), abs(rmssd_vs_hr.error))
}

-- Stage 5: Neural model binding (HuggingFace Hub)
model stress_detector {
  hub   "bitspark/autonomic-stress-classifier"
  task  classification

  input {
    dhr_autonomic       from pchr
    rt_ratio            from motor
    mean_iki            from motor
    rc_mean, se         from rppg
  }

  output stress_label   -- "low" | "moderate" | "high"
  output stress_prob    -- confidence vector over 3 classes
}

-- Stage 6: Watch conditions (rule-based)
watch stress_onset {
  when  dhr_autonomic > 15.0   -- genuine neural drive elevated
  and   regime        is turbulent
  and   rt_ratio      < 0.3    -- spinal, not cortical motor pattern
  and   csci          > 0.70   -- coupling maintained (not artefact)
  emit  "stress_onset"  confidence 0.8
}

watch stress_thermal_confound {
  when  dhr_metabolic > 8.0    -- thermal drive dominant
  and   dhr_autonomic < 5.0    -- true autonomic drive is low
  and   regime        is coherent
  emit  "thermal_elevation_not_stress"  confidence 0.85
}

-- Stage 7: Alignment and explanation
explain stress_audit {
  brut_prediction   = first_match(stress_onset, stress_thermal_confound)
  model_prediction  = stress_detector.stress_label
  model_confidence  = stress_detector.stress_prob

  match {
    agree when brut_prediction aligns model_prediction

    diverge when brut_prediction is "stress_onset"
            and model_prediction is "low" {
      -- Model says low stress; BRUT says onset. Expose the decomposition.
      trace pchr.dhr_autonomic
      trace pchr.dhr_metabolic
      trace skin_optics.t_skin
      trace motor.rt_ratio
      trace s_entropy.regime
      flag "Possible thermal confound: check t_skin and dhr_metabolic"
      flag "Check training regime: model may not separate PCHR components"
    }

    diverge when brut_prediction is "thermal_elevation_not_stress"
            and model_prediction is "high" {
      trace pchr.dhr_metabolic
      trace skin_optics.vasodilation
      trace skin_optics.t_skin
      flag "Model is classifying thermal tachycardia as stress"
      flag "dhr_autonomic is low; elevated HR is metabolic in origin"
    }
  }

  annotate model_prediction with {
    hr_components : pchr
    optical_state : skin_optics
    coupling      : coupling
    motor_state   : motor
  }
}
\end{lstlisting}

% =====================================================================
\section{Implementation Blueprint}
\label{sec:implementation}
% =====================================================================

\subsection{Compiler Architecture}

The \BS{} compiler is a four-stage pipeline implemented in TypeScript
and bundled by Vite into the BRUT Observatory web application.  The
four stages are:

\begin{enumerate}[leftmargin=*,nosep]
\item \textbf{Lexer.} Converts source text to a flat token stream.
  Token types: \texttt{KEYWORD}, \texttt{IDENT}, \texttt{NUMBER},
  \texttt{STRING}, \texttt{PUNCT}, \texttt{COMMENT} (discarded),
  \texttt{EOF}.

\item \textbf{Parser.} Recursive-descent parser over the grammar of
  \cref{sec:grammar}.  Produces an Abstract Syntax Tree (AST) whose
  node types correspond one-to-one to the block and statement types.
  Every AST node carries a \texttt{sourceLocation} field (line, column).

\item \textbf{Compiler.} Single-pass over the AST.  Performs
  signal-closure checking (\cref{thm:closure}), type checking
  (\cref{thm:type}), and emits a \emph{wiring plan}: a serialisable
  JSON object describing the evaluation order, the signal names,
  and the inversion functions to call at each step.

\item \textbf{Runtime.} Executes the wiring plan against the live
  signal bus provided by the Observatory's frame loop.  At each
  1\,Hz tick, the runtime evaluates all blocks in order, emits
  NDJSON trace entries to the sandbox output panel, and resolves
  model inference calls via the HuggingFace Inference API.
\end{enumerate}

\subsection{AST Node Types}

\cref{tab:ast} lists the TypeScript interface for each AST node type.

\begin{table}[H]
\centering
\caption{TypeScript AST node interfaces.}
\label{tab:ast}
\footnotesize
\begin{tabularx}{\columnwidth}{lX}
\toprule
\textbf{Interface} & \textbf{Key fields} \\
\midrule
\texttt{SourceBlock}   & \texttt{name}, \texttt{signals: string[]}, \texttt{rateHz} \\
\texttt{LayerBlock}    & \texttt{name}, \texttt{from: string[]}, \texttt{stmts: (InvertStmt | DeriveStmt)[]} \\
\texttt{InvertStmt}    & \texttt{output}, \texttt{inputSignal}, \texttt{given?: string[]}, \texttt{using}, \texttt{baseline?}, \texttt{sqrtCompress?} \\
\texttt{DeriveStmt}    & \texttt{output}, \texttt{expr: Expr}, \texttt{clamp?: [number,number]} \\
\texttt{DecomposeBlock}& \texttt{name}, \texttt{from: string[]}, \texttt{stmts: DecomposeStmt[]} \\
\texttt{TermStmt}      & \texttt{output}, \texttt{expr: Expr} \\
\texttt{CoordStmt}     & \texttt{output}, \texttt{expr: Expr} \\
\texttt{RegimeStmt}    & \texttt{input}, \texttt{cases: \{label, cond\}[]}, \texttt{otherwise} \\
\texttt{FailureStmt}   & \texttt{cases: \{label, cond\}[]}, \texttt{none} \\
\texttt{PairStmt}      & \texttt{name}, \texttt{predicted: Expr}, \texttt{observed: Expr} \\
\texttt{WatchBlock}    & \texttt{name}, \texttt{conds: Condition[]}, \texttt{emit: string}, \texttt{confidence} \\
\texttt{ModelBlock}    & \texttt{name}, \texttt{hub}, \texttt{task}, \texttt{inputs: InputBinding[]}, \texttt{outputs: string[]} \\
\texttt{ExplainBlock}  & \texttt{name}, \texttt{brutPred: Expr}, \texttt{modelPred: Expr}, \texttt{agree: Condition}, \texttt{divergeBranches: DivergeBranch[]}, \texttt{annotate?} \\
\bottomrule
\end{tabularx}
\end{table}

\subsection{Signal Bus Interface}

The runtime connects to the Observatory frame loop via a typed signal
bus interface:

\begin{lstlisting}[
  language={},
  numbers=none,
  basicstyle=\scriptsize\ttfamily,
  caption={TypeScript signal bus interface.},
  label={lst:bus}
]
interface SignalBus {
  get(name: string): number | string;
  subscribe(name: string,
    cb: (v: number|string, t: number) => void): void;
  emit(event: string, payload: TraceEntry): void;
}

interface TraceEntry {
  t: number;       // session ms
  block: string;
  step: string;
  line: number;
  in?: Record<string, number | string>;
  out?: Record<string, number | string>;
  eval?: boolean;
  value?: number | string;
  emit?: string;
  confidence?: number;
  flag?: string;
  traces?: Record<string, number | string>;
}
\end{lstlisting}

\subsection{HuggingFace Model Resolution}

Model blocks resolve their \texttt{hub} identifier against the
HuggingFace Inference API (\texttt{api-inference.huggingface.co}).
The runtime sends a POST request with the assembled input tensor
and receives a JSON response.  The task type (\texttt{classification},
\texttt{regression}, \texttt{token-classification}) determines the
response schema.  The runtime maps response fields to the declared
\texttt{output} identifiers by position.

For offline or low-latency use, the \texttt{hub} field may also
resolve to a local ONNX or TensorFlow.js checkpoint loaded into the
browser via \texttt{@xenova/transformers} or \texttt{onnxruntime-web}.
The compiler generates different call sites for remote vs local
resolution based on URL scheme: \texttt{"org/name"} for remote,
\texttt{"file:..."} for local.

% =====================================================================
\section{Discussion}
\label{sec:discussion}
% =====================================================================

\subsection{Relationship to Existing Explainability Methods}

\BS{} is not a post-hoc explainability method.  SHAP~\cite{lundberg2017shap}
and LIME~\cite{ribeiro2016lime} perturb inputs and measure output
changes, but they are agnostic to the physical meaning of the perturbed
features.  A SHAP value on the green channel pixel mean says nothing
about whether that pixel was carrying cardiac, vasodilatory, or melanin
information.  \BS{} is a \emph{pre-hoc} specification: the developer
writes what the physics says is happening, compiles it, and the runtime
shows whether the neural model agrees at each step.

The closest related paradigm is Concept Bottleneck
Models~\cite{koh2020concept}, which impose interpretable intermediate
representations.  \BS{} differs in that it does not modify the model
architecture --- the model remains a black box, and \BS{} provides an
external parallel pipeline that can be compared against it.  This makes
\BS{} applicable to any pre-trained model without retraining.

\subsection{The PCHR Decomposition as a Disambiguation Primitive}

The central insight of the sensor disambiguation
framework~\cite{sachikonye2024disambiguation} is that multi-sensor
data is a set of coupled projections from a five-dimensional partition
manifold.  A neural model that takes raw features as input does not
respect this manifold structure --- it may confound the thermal and
autonomic components of heart rate because they are correlated in
typical training data.  The \texttt{dhr\_autonomic} term in a \BS{}
script is the quantity that isolates genuine neural drive from thermal
and hypoxic contributions.  When the model disagrees with the
\BS{} trace on this term, it indicates a specific failure mode: the
model has not learned to disambiguate physiological state.

\subsection{Motor Channels as Independent Audit Signals}

The rambling--trembling decomposition~\cite{sachikonye2024posture}
provides a motor channel that is physiologically independent of the
optical rPPG channel.  Supraspinal rambling power reflects voluntary
cortical target-tracking; spinal trembling power reflects reflex
stabilisation.  A high-stress state should elevate trembling power
and lower the $\Ra/(\Ra+\Tr)$ ratio independently of cardiac
coherence.  A \BS{} \texttt{watch} block that requires both
$\Delta\HR_{\mathrm{auto}} > 15$ and $\Ra/(\Ra+\Tr) < 0.3$ is
thus requiring convergent evidence from two independent physiological
pathways, making the audit trail more robust to single-channel noise.

\subsection{Limitations and Future Work}

\BS{} in its current form has three limitations.

First, the language is synchronous: all blocks in a script are
evaluated at the same tick.  Future work should introduce
\texttt{async} blocks for model inference that can fire on a slower
clock than the main sensor loop.

Second, the \texttt{baseline} built-in requires a rolling window of
historical data.  A newly started session will not have a 7-day
baseline for $\HR_{\mathrm{int}}$.  A calibration block type should
be added to allow scripts to specify fallback constants.

Third, the diverge branch of \texttt{explain} currently emits a flag
string.  A more powerful design would allow \texttt{diverge} branches
to \emph{spawn} a secondary model query with the PCHR decomposition
terms as additional context features --- a structured prompting
pattern for language model-based interpretation.

% =====================================================================
\section{Conclusion}
\label{sec:conclusion}
% =====================================================================

We have presented \BS{}, a domain-specific language whose vocabulary
is the partition-theoretic quantity set of the Bounded Phase Space
framework.  The language provides six block types --- \texttt{source},
\texttt{layer}, \texttt{decompose}, \texttt{watch}, \texttt{model},
and \texttt{explain} --- whose semantics are formally defined, whose
correctness is governed by four theorems, and whose compiled output
is a structured JSON execution trace in which every intermediate
value is named, time-stamped, and attributed to the source line that
produced it.

The core contribution is that \BS{} makes the black-box problem in
physiological AI tractable without retraining: a developer writes a
physics-grounded script, binds a HuggingFace model as a named
participant, and the runtime shows exactly where the rule-based
trace and the neural prediction agree or diverge, with a pointer to
the specific decomposition term responsible.  The language is
implementable as a TypeScript compiler in the BRUT Observatory web
application and is designed to run in real time against live WebGPU
rPPG data.

\bibliographystyle{unsrt}
\bibliography{references}

\end{document}
